\textbf{COLUNI 2020 -- Questão 25}\par\medskip Leia o fragmento da reportagem com informações da Agência Brasil, veiculada pelo Portal FolhaPE em 22 de março de 2019.

\par

\textbf{Lama de Brumadinho chega ao Rio São Francisco}

\par

Segundo a Fundação SOS Mata Atlântica, as concentrações de ferro, manganês, cromo e cobre também estavam acima dos limites máximos permitidos pela legislação, o que evidencia o impacto da pluma de rejeitos de minério sobre o Alto São Francisco. Disponível em: <https://www.folhape.com.br/noticias/noticias/brumadinho/2019/03/22/NWS,99707,70,1370,NOTICIAS,2190-LAMA-BRUMADINHO-CHEGA-RIO-SAO-FRANCISCO.aspx>. Acesso em: 11 set. 2019. (Adaptado)

\par

O fragmento da reportagem acima informa sobre elementos químicos presentes na água do Rio São Francisco. Muitas das espécies químicas presentes em água estão na forma de íons, liberados a partir de compostos dissolvidos em água. Esses íons podem voltar a recombinar-se para dar origem ao composto original.

\par

Observe a tabela abaixo, na qual as espécies de alguns elementos estão representadas genericamente pelas letras A, B, C, D, E:

\par

\begin{center}
\begin{tabular}{|c|c|c|c|}
\hline
Espécie & Número de prótons & Número de elétrons & Número de nêutrons \\
\hline
A & 11 & 10 & 12 \\
\hline
B & 26 & 24 & 30 \\
\hline
C & 17 & 18 & 18 \\
\hline
D & 16 & 18 & 16 \\
\hline
E & 80 & 80 & 122 \\
\hline
\end{tabular}
\end{center}

\par

O composto formado por pares dessas espécies, constituído por cátions com carga +1 e ânions com carga -1, quando colocado em água, apresentará grande quantidade de íons separados (íons livres). Já outro composto formado por pares dessas espécies, constituído por cátions com carga +2 e ânions com carga -2, quando colocado em água, originará uma menor quantidade de íons livres.

\par

Com base nas informações fornecidas e nos conhecimentos sobre ciências, analise as seguintes afirmativas:

\par

I.   O composto formado pelas espécies A e C, quando colocado em água, apresentará grande quantidade de íons separados (íons livres). II. As espécies B, C e E correspondem, respectivamente, a cátion, ânion e átomo. III. Um dos elementos químicos citados no fragmento da reportagem está representado entre as espécies da tabela.

\par

Está CORRETO o que se afirma em:\begin{enumerate}
\item I e II, apenas.
\item I e III, apenas.
\item II e III, apenas.
\item I, II e III.\end{enumerate}
